\documentclass{amsdtx}
\usepackage[margin=.75in]{geometry}
\usepackage{graphicx}
\usepackage{float}
\usepackage{eqnarray,amsmath}
\usepackage{amssymb}
\usepackage{url}
\usepackage[dvipsnames]{xcolor}
\usepackage{cancel}
\usepackage{multicol}
\usepackage[Glenn]{fncychap}
\usepackage{esint}
%\usepackage{hyperref}
\definecolor{b}{HTML}{3C78D8}
\definecolor{t}{HTML}{134F5C}
\definecolor{g}{HTML}{38761D}
\definecolor{r}{HTML}{AB0000}
\definecolor{p}{HTML}{A64D79}
\usepackage[Glenn]{fncychap}
\usepackage{calligra}
\newcommand{\bs}[1]{\boldsymbol{#1}}

\usepackage{tipa}
\DeclareMathAlphabet{\mathcalligra}{T1}{calligra}{m}{n}
\DeclareFontShape{T1}{calligra}{m}{n}{<->s*[2.2]callig15}{}
\newcommand{\scriptr}{\mathcalligra{r}\,}
\newcommand{\boldscriptr}{\pmb{\mathcalligra{r}}\,}
%\usepackage{hyperref}



\usepackage{bm}

\newcommand*{\mline}[1]{%
\begingroup
    \renewcommand*{\arraystretch}{1.1}%
   \begin{tabular}[c]{@{}>{\arraybackslash}p{2cm}@{}}#1\end{tabular}%
  \endgroup
}


\usepackage{listings}

\definecolor{codegreen}{HTML}{38761D}
\definecolor{codeblue}{HTML}{000076}
\definecolor{codegray}{rgb}{0.2,0.2,0.2}
\definecolor{codepurple}{HTML}{AB0000}
\definecolor{backcolour}{rgb}{0.95,0.95,0.95}

\lstdefinestyle{mystyle}{
    backgroundcolor=\color{white},   
    commentstyle=\color{codegreen},
    keywordstyle=\color{codepurple},
    numberstyle=\ttfamily\scriptsize\color{codeblue},
    stringstyle=\color{b},
    basicstyle=\ttfamily\footnotesize,
    breakatwhitespace=false,         
    breaklines=true,                 
    captionpos=b,                    
    keepspaces=true,                 
    numbers=left,                    
    numbersep=5pt,                  
    showspaces=false,                
    showstringspaces=false,
    showtabs=false,                  
    tabsize=4
}
\lstset{style=mystyle}



\newcolumntype{M}[1]{>{\centering\arraybackslash}m{#1}}
\newcolumntype{N}{@{}m{0pt}@{}}

\newcommand{\vfs}[1]{\ensuremath \textbf{\textit{f}}\kern-0.1cm_{s_{#1}}}
\newcommand{\fs}[1]{\ensuremath f\kern-0.07cm_{s_{#1}}}
\newcommand{\F}{\ensuremath \textbf{F}\kern-0.05cm}
\newcommand{\cc}{\ensuremath\mu_{cc}}
\newcommand{\cg}{\ensuremath\mu_{cg}}
\newcommand{\e}{\ensuremath{\textbf{{\bf e}}}}
\newcommand{\M}{\ensuremath{\textbf{M}}}
\newcommand{\R}{\ensuremath{\textbf{r}}}
\newcommand{\U}{\ensuremath{\textbf{u}}}
\newcommand{\st}{\ensuremath{\textbf{\bf t}}}
\newcommand{\fnf}{\ensuremath{F_{\rm N_{\it f}}}}
\newcommand{\Fnf}{\ensuremath{\text{$\F_{\rm N_{\it f}}$}}}
\newcommand{\dvec}[1]{\ensuremath\overset{\text{\tiny$\bm\leftrightarrow$}}{#1}}

\usepackage{enumitem}

\title{\sc OpenRocket-Based Rocket Control Simulation}
\author{\sc Marc D Nichitiu~|~MIT Rocket Team Simulations}
\date{\sc Fall 2025}
\begin{document}
\maketitle    
\begin{center}
	\it In this document, we outline the purpose, structure, and uses of the Java-based OpenRocket software in scripting control simulation for amateur rocketry purposes.
\end{center}
\begin{multicols}{2}
	\tableofcontents
\end{multicols}

\section{Introduction}
OpenRocket is a very versatile open-source software written in Java, allowing for multi-platform interoperability. OpenRocket is quite flexible and reliable for subsonic regimes, with moderately accurate of supersonic simulation. For amateur rocketry, it is well within the accuracy of the manufacturing modeled. 

Using the Python {\tt jpype} module, it is possible to instantiate Java objects, thus scripting the OpenRocket simulation routines. It is additionally possible to use Simulation Listeners to actuate rocket components in flight as a function of flight data.

We thus seek to apply these in order to optimize our rocket designs.
\section{OpenRocket}
\subsection{Structure}
\subsection{Components}
\subsection{Simulation Flow}
\section{Roll Tabs}
\subsection{Design}
\subsection{Implementation}
\subsection{Control Algorithm}
\subsection{Simulation}
\section{Airbrakes}
\subsection{Design}
\subsection{Implementation}
\subsection{Control Algorithm}
\subsection{Simulation}



\end{document}














































